\section{Data, Provenance and Experimental Design}
\label{sec:data-design}

\subsection{Data universe and dates}

The legacy corpus contains different evidence domains: long-run gold and silver
time series, LSE (London Strategic Edge) market and option inputs, mine-level
operating data and synthetic DCF examples. Each retains its own provider,
frequency, date range, unit and measure. The operating domain is quarterly and
mine-level: Team~4 estimates on 2019Q1--2025Q4 Barrick reports, with shorter
effective samples where leading observations are missing; 2026Q1--Q2 enter
as actuals used only for level checks, and forecasts cover 2026Q3--2030Q4, of
which the valuation counts only flows after 2 September 2026
(Section~\ref{sec:company-perimeter}). The 2 April 2026 GLD option snapshot
and standalone Team~4 GBM paths are a separate legacy domain, excluded from
the integrated valuation \cite{team4}. Gold histories, option-calibration
surfaces and operating panels are not pooled.

The copper extension is another declared experiment: historical HG/HXE
bid/ask surfaces, contemporaneous matching futures and date-specific Treasury
curves are distinct from both GLD and the mine panel. Its preliminary temporal
pilot samples ten dates across 9 July--2 September 2026, including the final
three consecutive dates, with approximately 100 spaced real locations per
date. It is not a reproduction of the 31-date gold OOS panel. Missing, stale
or wide-spread quotes remain exclusions; no synthetic quotes complete a date.
The surviving contract catalogue prevents a complete point-in-time universe
and expired-option history is unavailable from this API
\cite{ib-historical-limitations}. Accepted dates and actual forecast gaps are
reported in Section~\ref{sec:copper-extension}.

\subsection{Provenance and reproducibility}

Every quantitative result included in the final study must be traceable through an input hash, as-of date, configuration, seed where applicable, executable stage, output path and tolerance. Exact reuse identifies material preserved from a source article; edited reuse identifies rearrangement or notation harmonisation; recomputed identifies a result created by a fixed, versioned executable run. A gap never authorises a fabricated value.

\subsection{Common notation, units and model interfaces}

The notation distinguishes the physical measure $\mathbb P$ from the pricing
measure $\mathbb Q$, GLD in USD/share from gold in USD/oz, annualised decimals
from percentage points, and enterprise from equity value. Production is
$\mathrm{Prod}$ or italic $Q$ in Chapter~\ref{ch:dcf}, distinct from $\mathbb Q$.
Ore is in tonnes, grade in g/t, recovery decimal, and production in troy
ounces of $31.1035$ grams, usually reported in thousands. Gold CoS is USD
per attributable ounce sold. Copper prices and CoS are USD/lb, attributable
sales are kt, and $1~\mathrm{t}=2204.6226218488~\mathrm{lb}$. Values use an
explicit scale, normally USD millions; per-share values require a reconciled
diluted-share denominator. Already attributable volumes are never scaled by
ownership twice.

Several standard symbols keep their conventional, local meaning in different
chapters. In Chapter~\ref{ch:econometrics}, $\theta_j$ are moving-average
coefficients, $p$ and $q$ are model orders and $L$ is the lag operator
($Ly_t=y_{t-1}$); in the GARCH models of Chapters~\ref{ch:econometrics}
and~\ref{ch:gold-econometrics}, $\alpha$ and $\beta$ are the ARCH and GARCH
coefficients. In Chapters~\ref{ch:sv-jumps} and~\ref{ch:calibration}, $\kappa$
and $\theta$ (written $\theta_v$ in Chapter~\ref{ch:calibration}) are the
mean-reversion speed and long-run level of the variance, $\alpha$ also denotes
the Carr--Madan damping coefficient, and in the Hawkes intensity $\alpha$ and
$\beta$ are the excitation size and decay rate; the branching ratio is
$\eta=\alpha/\beta$, distinct from event counts $n$, and $A,B$ (or $C,D$) are
the Riccati coefficients of affine characteristic functions. In
Chapter~\ref{ch:dcf}, $\kappa$ is the mean-reversion speed of the WACC process.

\subsection{Evaluation design}

Model fit, predictive performance, numerical convergence and economic sensitivity are separate questions. In-sample residual improvement is not evidence of out-of-sample superiority. Likewise, source-model parity is a prerequisite for integration, not a new economic finding.

Error metrics must use the scale relevant to the valuation. Team~4's
12-quarter cost holdout against a random walk reports RMSE $0.0504$ versus
$0.1512$ and MAPE $0.603\%$ versus $1.77\%$ on the log master chain
\cite{team4}. A log error of $0.05$ is roughly a 5\% level error; MAPE on
log values is not MAPE in USD/oz and depends on the log level and unit scale.
The RMSE ratio is the useful comparison, not the stated MAPE level.

Coverage and uncertainty sources must be explicit. The 95\% cost bands and
80\% production bands describe operating uncertainty; the Chapter~\ref{ch:unified-experiment}
quantiles vary gold prices and discount rates with production and costs fixed.
The copper extension adds copper-price uncertainty under declared dependence
scenarios. These quantities are neither interchangeable nor complete
corporate confidence intervals.

Copper reserved-strike scores test interpolation. Temporal scores fit only
origin-date quotes and project variance/intensity to the next accepted date;
target forward and as-of rates condition repricing. Common target support,
origin-mean and inside-hull persistence benchmarks, date-equal and pooled
losses are retained separately. Conditional IV repricing does not test a
physical price forecast or the accounting cash-flow bridge.

\paragraph{Required physical forecasting and validation.}
The corporate DCF requires commodity paths under $\mathbb P$, followed by revenues and reconciled FCFF discounted at a compatible WACC. The four model families can be specified under both measures; option calibration alone does not identify their physical parameters. The next estimation stage must jointly constrain historical underlying dynamics and option prices through explicit risk premia. Diffusion coefficients and instantaneous correlations must remain compatible with an equivalent change of measure; variance drifts and jump compensators can change. For Hawkes jumps, preservation of the chosen intensity family must be demonstrated rather than assumed.

Three validation targets are required for both metals. Option repricing evaluates a $\mathbb Q$ pricing function on later observed surfaces, conditional on the stated target-date inputs; observed surface evolution itself occurs under $\mathbb P$. Separate origin-only tests must score physical commodity prices and predictive distributions against realised outcomes and historical benchmarks. Finally, Barrick sales, costs, revenues and FCFF forecasts must be compared with subsequently published reports on the same perimeter. Neither IV persistence scores nor comparison of a DCF with the stock quote substitutes for these physical tests. The present study has not completed either physical validation stage.

\subsection{Research boundaries}

This document is a reproducible research compilation and contains no
investment recommendation. Not every element of the eight source studies is
reproduced. The theoretical Chapters~\ref{ch:probability}--\ref{ch:monte-carlo}
and~\ref{ch:sv-jumps} retain the mathematics but omit the source studies'
code listings, synthetic illustrations and their numerical outputs, so the
variance-reduction benchmark of Chapter~\ref{ch:monte-carlo} is described by its
design rather than its results; the mine-level production plots of the Team~4
study are not repeated in Chapter~\ref{ch:operations}; and several
source-study diagnostic figures for the historical 12 August calibration in
Chapter~\ref{ch:calibration} are replaced by their regenerated counterparts for
the current 2 September 2026 snapshot. These omissions concern illustrations,
not inputs to the valuation.
